\chapter{Conformação por convexidade}

\eyebrow{forma e musculatura, a partir do perfil}

\vspace{4pt}
A conformação descreve a forma e o desenvolvimento muscular da carcaça. A
norma de avaliação a define pela \textbf{convexidade do perfil}: musculatura plena
empurra o contorno para fora (convexo); musculatura pobre, com osso aparente, o afunda
para dentro (côncavo). O Apex transcreve exatamente essa regra em uma medida.

\section{O método}

A medida de \textbf{convexidade integral-invariante} (Manay et al.\ 2004; Pottmann
et al.\ 2009) calcula, em cada ponto do contorno da carcaça, quanto de um pequeno
disco cai dentro da silhueta. Essa área carrega a curvatura com sinal sem
diferenciar --- por isso é robusta ao ruído da máscara --- e produz um valor:

\begin{center}
\ttfamily $c > 0$ \rm convexo (musculoso) \quad\textbullet\quad
\ttfamily $c = 0$ \rm reto \quad\textbullet\quad
\ttfamily $c < 0$ \rm côncavo (osso aparente)
\end{center}

Ela é \textbf{analítica} (sem treino, sem pesos), \textbf{invariante} a
rotação, translação, escala e iluminação, e usa apenas a silhueta segmentada
--- sem cor nem textura. O valor é promediado por região anatômica: \emph{leg}
(perna), \emph{loin} (lombo) e \emph{shoulder} (paleta), com a pose pendurada
(jarrete no topo) fixando o eixo.

\section{Lendo o mapa de convexidade}

A aba \textbf{Conformation} (Figura~\ref{fig:analysis-conformation}) mostra, por
carcaça, o contorno colorido pela convexidade sobre uma foto esmaecida: \textbf{azul onde
convexo} (musculoso), \textbf{vermelho onde côncavo}. A convexidade média de cada região é
anotada, e abaixo do mapa os índices de perna/lombo/paleta são listados com uma
pílula de grau estimado.

\shot{analysis-conformation.jpeg}{Analysis $\rightarrow$ Conformation. O mapa de convexidade
reproduz a anatomia que a norma de avaliação descreve: as pernas aparecem em azul
(convexas, musculosas), o recorte medial entre elas aparece em vermelho (uma reentrância). Índices
por região e um grau estimado aparecem abaixo de cada mapa.}{analysis-conformation}

Repare como o mapa é consistente com a anatomia sem qualquer rótulo: as pernas são
azuis, o recorte entre elas é vermelho, o lombo é quase neutro. Essa
correspondência é verificável a olho.

\begin{caution}[title=O grau estimado não é validado]
O mapa de convexidade e os índices por região são \textbf{medidas objetivas}.
O \emph{grau estimado} derivado deles (mostrado com uma pílula \tele{est.} em
âmbar) \textbf{não é validado} --- o experimento oráculo do projeto mostrou que
a conformação não é recuperável de forma confiável a partir de uma única imagem 2D neste tamanho de
amostra (Estação~13). Use o grau como referência indicativa, a ser validado à
medida que o conjunto de dados pareado cresce.
\end{caution}

\section{Onde os números aparecem}

A conformação roda automaticamente durante a Análise (quando o fundo é removido).
Para cada carcaça analisada, o Inspector mostra as convexidades de perna / lombo / paleta,
o índice consolidado e o grau estimado
(Figura~\ref{fig:carcass-images}, à direita). A exportação em CSV inclui todos eles.
